% SPDX-License-Identifier: GPL-3.0-or-later
% Copyright (C) 2026 Alessandro Greco
\chapter{Измерение сети: Tor Metrics и OONI}\label{cap:24-metrics}
\aggiornatoal{23 сентября 2026 г.}{данные Tor Metrics по 20 сентября 2026 г.; консенсус от 23 сентября 2026 г., 15:00; спецификации OONI ts-023 и ts-030}

\begin{obiettivi}
К концу этой главы вы сможете:
\begin{itemize}[nosep]
  \item описывать данные, публикуемые Tor Project, и источники этих данных;
  \item объяснять, как оценивается число пользователей в сети, которая специально не должна знать своих пользователей, и насколько неопределённы такие оценки;
  \item извлекать сетевую статистику из реального консенсуса (Лабораторная работа~15);
  \item объяснять, как OONI измеряет блокировку Tor с точки зрения пользователей.
\end{itemize}
\end{obiettivi}

Сеть анонимности создаёт для операторов необычную проблему: для улучшения сети её нужно измерять, но измерение означает сбор информации именно о том, что она призвана защищать. Подход Tor Project — измерять минимально необходимое, агрегировать и публиковать результаты и документировать каждый шаг так, чтобы расчёты мог воспроизвести любой исследователь~\cite{metricsreproducible}. В этой главе мы проследим путь данных от источников до графиков, а затем посмотрим на сеть снаружи через измерения OONI.

% ---------------------------------------------------------------------------
\section{Источники данных}\label{sec:metrics-fonti}

Почти всё, что публично известно о сети Tor, происходит из документов, уже встречавшихся нам:
\begin{description}
  \item[Консенсус и дескрипторы] (\cref{cap:12-directory}): список ретрансляторов, флаги, весовые коэффициенты пропускной способности, версии и политики выхода.
  \item[Дополнительные дескрипторы (extra-info)] Публикуются каждым ретранслятором вместе с основным дескриптором и содержат историю объёма переданных байтов и агрегированную статистику, включая число запросов к каталогу по странам.
  \item[Активные измерения] OnionPerf периодически загружает файлы через Tor из выбранных точек по всему миру, записывая время передачи и ошибки.
\end{description}
\textbf{CollecTor} архивирует эти документы, начиная с ранних лет сети, и предоставляет их для загрузки. \textbf{Tor Metrics} строит на их основе графики и таблицы, сопровождая их CSV-данными и документированными методиками. \textbf{Relay Search} позволяет находить сведения об отдельных ретрансляторах. Консенсус, использованный в лабораторных работах этой книги (\file{fonti/collector}), был получен из CollecTor.

% ---------------------------------------------------------------------------
\section{Подсчёт анонимных пользователей}\label{sec:metrics-utenti}

Tor не регистрирует пользователей, и ни один ретранслятор не может определить, принадлежат ли два соединения одному человеку. Как тогда сеть оценивает, что ею пользуются примерно три миллиона клиентов? Метод косвенный (\cref{fig:stima}). Каждый клиент периодически загружает консенсус, а запросы к каталогу можно подсчитывать, не идентифицируя запрашивающих. Ретрансляторы, обслуживающие запросы к каталогу, ежедневно считают успешные запросы, группируя их по стране, определённой по адресу источника. Они публикуют агрегированные значения в дескрипторах extra-info, округляя числа по странам вверх до кратных восьми, чтобы маленькие точные значения раскрывали меньше информации~\spec{dir-spec}{Extra-info document format}.

\begin{figure}[htbp]
\centering
% SPDX-License-Identifier: GPL-3.0-or-later
% Copyright (C) 2026 Alessandro Greco
%% Альтернативное описание EPUB.
\altfig{Косвенная дневная оценка Tor Metrics с отдельными рядами для прямых клиентов реле и клиентов мостов. Реле каталогов и мосты подсчитывают успешные запросы документов каталога примерно за 24 часа. В extra-info поле dirreq-v3-reqs числа по странам округляются вверх до кратных 8. CollecTor архивирует эти дескрипторы. Tor Metrics вычитает 4 из допустимых округлённых значений, оценивает число наблюдаемых запросов r(R) и долю покрытия frac, затем вычисляет floor(r(R)/(frac умножить на 10)). Число 10 — предположение о среднем числе запросов на одного клиента в день. Дни с frac ниже 10 процентов или выше 110 процентов исключаются. Итог означает среднее число одновременно подключённых клиентов, а не отдельных людей.}
% Косвенная оценка среднего числа клиентов по агрегированным запросам
\begin{tikzpicture}[font=\sffamily\scriptsize,node distance=0.5cm,
  box/.style={draw,rounded corners=2pt,align=center,minimum height=0.9cm,text width=2.3cm}]
  \node[box,fill=torviola!10] (r) {реле и мосты\\считают запросы\\каталога};
  \node[box,right=of r] (x) {extra-info\\\texttt{dirreq-v3-reqs}};
  \node[box,right=of x] (c) {CollecTor\\архивирует};
  \node[box,right=of c,fill=torverde!15] (m) {Tor Metrics:\\$\lfloor r(R)/(\mathit{frac}\cdot10)\rfloor$};
  \draw[->] (r) -- (x); \draw[->] (x) -- (c); \draw[->] (c) -- (m);
  \node[below=2pt of r,font=\sffamily\tiny,text=torgrigio,align=center] {период около\\24 часов};
  \node[below=2pt of x,font=\sffamily\tiny,text=torgrigio,align=center] {по странам, кратно\\8 (затем $-4$)};
  \node[below=2pt of c,font=\sffamily\tiny,text=torgrigio,align=center] {$\mathit{frac}$: оценка\\покрытия};
  \node[below=2pt of m,font=\sffamily\tiny,text=torgrigio,align=center] {среднее число клиентов\\не число людей};
  \node[below=0.95cm of c,font=\sffamily\tiny,text=avviso,align=center] {день пропускают при $\mathit{frac}<10\%$ или $\mathit{frac}>110\%$};
\end{tikzpicture}

\caption{От успешных запросов каталога к косвенной оценке \emph{среднего числа подключённых клиентов}, а не отдельных людей. Реле каталогов и мосты дают отдельные агрегированные ряды в дескрипторах extra-info; количества запросов по странам в \texttt{dirreq-v3-reqs} округляются вверх до кратных 8. CollecTor архивирует данные; Tor Metrics вычитает 4 из допустимых значений, оценивает долю охвата $\mathit{frac}$ и применяет формулу $\lfloor r(R)/(\mathit{frac}\cdot10)\rfloor$. Число 10 — предположение о среднем количестве запросов клиента в день. Дни при $\mathit{frac}<10\%$ или $\mathit{frac}>110\%$ исключаются.}
\label{fig:stima}
\end{figure}

Tor Metrics собирает эти значения и выполняет три шага~\cite{metricsreproducible}:
\begin{enumerate}
  \item Компенсирует округление, вычитая половину интервала округления из каждого значения.
  \item Оценивает \textbf{долю}, $\mathit{frac}$, всех запросов к каталогу сети, покрываемую ретрансляторами, предоставившими статистику, с учётом времени их работы и объёма обслуженных запросов.
  \item Делит зарегистрированные запросы на эту долю, получая оценку общего числа запросов в сети, затем делит на 10:
\[
r(N) = \left\lfloor \frac{r(R)}{\mathit{frac} \cdot 10} \right\rfloor.
\]
\end{enumerate}
Делитель 10 — предположение. Непрерывно работающий клиент делает примерно 15 запросов в сутки, но многие клиенты подключены лишь часть дня, поэтому Tor Project выбрал 10 как среднее значение. Оценки с долей охвата ниже 10\,\% или выше 110\,\% отбрасываются.

\begin{attenzione}[Что означает \enquote{пользователи}]
Результат оценивает число \emph{в среднем подключённых клиентов}, а не отдельных людей. Клиент, работающий два часа в день, обычно выполняет меньше запросов к каталогу, чем клиент, работающий круглосуточно, но фиксированного соотношения между ними нет; несколько программ с Tor на одной машине считаются несколькими клиентами. Изменения поведения ПО, например менее частая загрузка консенсуса, меняют оценку даже при неизменном числе людей. Поэтому графики полезнее для анализа \emph{изменений}, чем абсолютных значений.
\end{attenzione}

\paragraph{Значения на эталонную дату книги.} С 1 августа по 20 сентября 2026 года оценка ежедневно подключённых напрямую пользователей колебалась между 2,8 и 3,5 миллиона, при доле охвата от 57\,\% до 73\,\%. Среднее значение для Италии составляло примерно 60\,000, но ежедневные результаты имеют широкие ожидаемые интервалы: например, 20 сентября — примерно от 34\,000 до 111\,000. Изменения внутри таких диапазонов сами по себе не доказывают наличие события. В тот же день оценка числа пользователей мостов составляла около 98\,000 для obfs4, 60\,000 для Snowflake и 35\,000 для WebTunnel.

\subsection{Выявление цензуры}

Metrics использует те же временные ряды для выявления потенциальных событий цензуры. Для стран со значительным числом пользователей система сравнивает недельное изменение с типичными изменениями в других странах и выделяет даты, когда конкретная страна выходит за ожидаемый интервал~\cite{metricsreproducible}. Резкое падение числа прямых пользователей одновременно с ростом пользователей мостов — характерный признак блокировки. Это указание, а не доказательство. Для подтверждения нужны измерения внутри страны, и здесь используется OONI.

% ---------------------------------------------------------------------------
\section{Производительность}\label{sec:metrics-prestazioni}

OnionPerf из фиксированных точек измеряет время загрузки файлов разных размеров через Tor, ошибки и задержку цепи — интервал между отправкой HTTP-запроса и получением заголовков ответа~\cite{metricsreproducible}. 20 сентября 2026 года медианная задержка до публичного сервера составляла около 110--140\,мс из немецкой точки измерения, 280\,мс из США и более 500\,мс из Гонконга. Географическое расстояние до основного сосредоточения сети в Европе и Северной Америке имеет значение. Некоторые точки измерения используют Conflux — систему многопутевой передачи, описанную в \cref{cap:31-futuro}, что позволяет проводить сравнения при сходных условиях.

% ---------------------------------------------------------------------------
\section{Лабораторная работа: чтение реального консенсуса}\label{sec:lab-metrics}

\begin{laboratorio}[15]{Получение сетевой статистики из консенсуса}
\textbf{Цель.} Извлечь основные сетевые показатели из реального консенсуса CollecTor с помощью нового модуля \file{labs/ru/код/torlab/metrics.py}.\\
\textbf{Предварительные условия.} Лабораторная работа~6 (\file{directory.py}) и Лабораторная работа~7 (\file{path\_selection.py}); работающая сеть не требуется.\\
\textbf{Оценочное время.} 1 час.
\end{laboratorio}

Модуль добавляет функцию \texttt{summary}, повторно используя анализатор консенсуса из Лабораторной работы~6, а также две функции оценки пользователей. При равных значениях счётчики дополнительно сортируются по имени. Без этого их порядок менялся бы между запусками из-за порядка обхода множеств в Python. Чтобы воспроизвести данные исторического консенсуса от 23 сентября 2026 года, 15:00, сначала подготовьте файл из корня репозитория командой \texttt{make -C labs/ru data} (для загрузки нужен Интернет). Применительно к этому консенсусу (вывод сокращён):
\begin{lstlisting}[language=bash,numbers=none]
$ PYTHONPATH=labs/ru/код python3 -m torlab.metrics labs/ru/данные/collector/consensus-microdesc-2026-09-23-15
relay: 9358
флаги: {'Running': 9358, 'Valid': 9358, 'Fast': 8961, 'Stable': 8199,
       'V2Dir': 7674, 'Guard': 5306, 'HSDir': 4050, 'Exit': 3177,
       'BadExit': 32, 'MiddleOnly': 28, 'StaleDesc': 21, 'Authority': 10}
общая пропускная способность: 198927900
доля пропускной способности выходных узлов: 0.3308947161257923
доля пропускной способности сторожевых узлов: 0.7010271560701139
версии: [('Tor 0.4.9.12', 6027), ('Tor 0.4.9.11', 2520), ...]
сети /16: [('64.65', 771), ('192.42', 636), ('204.8', 598), ...]
\end{lstlisting}

Несколько наблюдений:
\begin{itemize}
  \item 9\,358 ретрансляторов согласуются с графиком Tor Metrics, который на 20 сентября показывает 9\,416 ретрансляторов и 2\,333 моста.
  \item Примерно треть ретрансляторов имеет флаг \texttt{Exit} и обеспечивает примерно треть измеренной пропускной способности; на сторожевые узлы приходится 70\,\%. Эти доли пересекаются, поскольку один ретранслятор может иметь оба флага. Дисбаланс компенсируется весовыми коэффициентами пропускной способности, обсуждаемыми в \cref{sec:percorso-pesi}.
  \item \texttt{BadExit} (32 ретранслятора) и \texttt{MiddleOnly} (28) отражают решения центров каталогов о подозрительных или неверно настроенных ретрансляторах.
  \item Десять ретрансляторов имеют флаг \texttt{Authority}, но в голосовании по консенсусу участвуют только девять центров каталогов (\cref{sec:arch-ruoli}). Десятый — \texttt{Serge}, центр каталогов мостов: флаг у него есть, но он не голосует.
  \item Большинство ретрансляторов использует современные версии 0.4.9.
  \item В нескольких сетях /16 размещены сотни ретрансляторов, обычно принадлежащих крупным операторам выходных узлов. Правило не более одного ретранслятора из одной сети /16 в цепи (\cref{sec:percorso-vincoli}) ограничивает зависимость от одной сети.
\end{itemize}
Сумма весов пропускной способности в консенсусе представляет единицы пропускной способности, измеренные центрами каталогов, а не фактически переданный трафик. Для оценки трафика Metrics использует историю переданных байтов из дескрипторов extra-info.

% ---------------------------------------------------------------------------
\section{OONI}\label{sec:metrics-ooni}

\emph{Open Observatory of Network Interference} (OONI) измеряет интернет-цензуру с точки зрения пользователей. Добровольцы по всему миру устанавливают OONI Probe, который периодически выполняет тесты и публикует открытые результаты. Некоторые тесты специально относятся к Tor~\cite{oonitor}:
\begin{description}
  \item[\texttt{tor}] Проверяет доступность выбранных центров каталогов и мостов. Тесты не ограничиваются установлением TCP-соединения и выполняют действия, специфичные для протокола, например рукопожатие obfs4.
  \item[\texttt{torsf}] Пытается запустить tor через Snowflake, записывая, завершилась ли первоначальная загрузка сети успешно и сколько времени он занял.
\end{description}
Спецификации тестов содержат защитные меры, предотвращающие публикацию приватных адресов мостов: их включение в результаты фактически передало бы эти адреса цензорам. Metrics и OONI дополняют друг друга. Metrics может показать падение числа пользователей из некоторой страны; OONI позволяет проверить, действительно ли ретрансляторы или мосты были недоступны изнутри этой страны, и определить технику блокировки.

\begin{approfondimento}[Измерять, не причиняя вреда]
Каждая статистика, публикуемая ретранслятором, является компромиссом. Чем выше детализация, тем полезнее данные, но тем выше риск раскрытия информации. Поэтому значения округляются и агрегируются по дням и странам. Исследования сети следуют тем же принципам (\cref{sec:rp-perche}): собирать минимум необходимого, агрегировать до публикации и не хранить лишние данные.
\end{approfondimento}

% ---------------------------------------------------------------------------
\section*{Итоги}
\begin{itemize}
  \item Сетевые данные происходят из консенсусов, дескрипторов и документов extra-info, архивируемых CollecTor, а также из активных измерений OnionPerf. Tor Metrics анализирует их по документированным и воспроизводимым методикам.
  \item Число пользователей оценивается по запросам к каталогу: зарегистрированные запросы делятся на покрываемую долю, затем на 10. Это оценка среднего числа подключённых клиентов, особенно полезная для анализа изменений.
  \item Те же данные могут указывать на возможные события цензуры, требующие подтверждения локальными измерениями.
  \item Реальный консенсус позволяет получить число relay, флаги, доли пропускной способности, версии и концентрацию по сетям.
  \item OONI измеряет доступность центров каталогов и мостов и работу Snowflake изнутри стран.
\end{itemize}

% ---------------------------------------------------------------------------
\section*{Упражнения}

\begin{esercizio}[1]\label{es:24-dieci}
Если все клиенты начнут загружать консенсус вдвое реже, что произойдёт с оценкой числа пользователей при неизменном числе людей? Как интерпретатор графиков мог бы заметить такое изменение?
\end{esercizio}

\begin{esercizio}[1]\label{es:24-frazione}
За один день ретрансляторы, предоставившие статистику, зарегистрировали двадцать миллионов запросов, а оценённая доля охвата составила 60\,\%. Сколько пользователей оценит Metrics? Что произойдёт, если доля охвата составит 8\,\%?
\end{esercizio}

\begin{esercizio}[2]\label{es:24-flag}
Расширьте \texttt{summary}, чтобы вычислять долю пропускной способности ретрансляторов, имеющих одновременно \texttt{Guard} и \texttt{Exit}, а также общую взвешенную вероятность того, что выходной узел цепи принадлежит одной из пяти наиболее представленных сетей /16.
\end{esercizio}

\begin{esercizio}[2]\label{es:24-censura}
Загрузите из Metrics временные ряды прямых пользователей и пользователей мостов для страны с документированным событием блокировки Tor. Опишите, что происходит с двумя рядами в даты блокировки.
\end{esercizio}

\begin{esercizio}[3]\label{es:24-privacy}
Статистика по странам округляется вверх до кратных восьми. Опишите атаку, которую это округление затрудняет, и объясните, почему оно не защищает от атакующего, контролирующего множество ретрансляторов.
\end{esercizio}

% ---------------------------------------------------------------------------
\section*{Дополнительная литература}
\begin{itemize}
  \item Страница Tor Project \titolo{Reproducible Metrics}~\cite{metricsreproducible}, объясняющая методы построения отдельных графиков.
  \item Спецификации тестов OONI~\cite{oonitor} и OONI Explorer: \url{https://explorer.ooni.org/}.
  \item Формат extra-info descriptor: \spec{dir-spec}{Extra-info document format}.
\end{itemize}
